\section{Универсальность математических моделей}
\label{sec:universality-models}

\textbf{Универсальность математической модели} означает, что одна и та же
математическая структура может описывать широкий класс объектов различной
физической природы после соответствующей интерпретации переменных и
параметров. Она возникает благодаря абстрагированию, математическим
аналогиям, безразмеризации и общим законам баланса.

\subsection{Абстрагирование}

При моделировании из реального объекта выделяют только характеристики,
существенные для поставленной цели. Например, разные тела при исследовании
поступательного движения могут быть представлены материальной точкой:
\[
m\ddot{\mathbf r}=\mathbf F.
\]
Чем меньше предметно-специфических деталей содержит модель, тем шире её
потенциальная область применения, но тем тщательнее нужно контролировать
потерянные эффекты.

\subsection{Математические аналогии}

Одно уравнение
\[
\dot x=-kx
\]
описывает радиоактивный распад, простейший разряд конденсатора, линейное
охлаждение и другие процессы. Уравнение
\[
\frac{\partial u}{\partial t}=D\Delta u
\]
возникает в теплопроводности, диффузии вещества и ряде вероятностных и
биологических моделей.

Совпадение математической формы позволяет переносить уже известные методы
анализа и численного решения. Однако математическая аналогия не означает
полной физической идентичности: различаются смысл переменных, допустимые
параметры, начальные и граничные условия и область применимости.

\subsection{Безразмеризация}

Пусть характерные масштабы задачи равны $L$, $T$, $U$. Вводят
\[
\mathbf x=L\widetilde{\mathbf x},
\qquad
t=T\widetilde t,
\qquad
u=U\widetilde u.
\]
После подстановки исходные размерные параметры объединяются в
безразмерные комплексы.

Например, для переноса и диффузии
\[
\frac{\partial c}{\partial t}
+\mathbf u\cdot\nabla c
=D\Delta c
\]
при масштабе времени $T=L/U$ получаем
\[
\frac{\partial c}{\partial\widetilde t}
+\widetilde{\mathbf u}\cdot\widetilde\nabla c
=\frac1{Pe}\widetilde\Delta c,
\qquad
\boxed{Pe=\frac{UL}{D}}.
\]
В гидродинамике аналогично возникает число Рейнольдса
\[
\boxed{Re=\frac{\rho UL}{\mu}}.
\]

\subsection[Пи-теорема Бэкингема]{$\Pi$-теорема Бэкингема}

Центральным результатом теории размерностей является $\Pi$-теорема.
Пусть физическая зависимость включает $n$ размерных величин, а ранг системы
их размерностей относительно основных единиц равен $r$. Тогда зависимость
можно переписать через
\[
\boxed{n-r}
\]
независимых безразмерных комплексов $\Pi_1,\ldots,\Pi_{n-r}$:
\[
F(\Pi_1,\ldots,\Pi_{n-r})=0.
\]
Это позволяет резко уменьшить число параметров эксперимента и выявить
определяющие режимы процесса.

Чтобы две системы были подобны, недостаточно совпадения одного критерия:
должны совпадать определяющие безразмерные параметры, безразмерная геометрия
и соответствующие начальные и граничные условия.

\subsection{Масштабные законы}

Многие универсальные зависимости имеют степенной вид
\[
Y\sim X^\alpha.
\]
Для геометрически подобных тел
\[
S\sim L^2,
\qquad
V\sim L^3,
\qquad
\frac SV\sim L^{-1}.
\]
Для диффузии характерное время оценивается как
\[
\boxed{T\sim\frac{L^2}{D}}.
\]
Такие оценки часто позволяют определить доминирующие механизмы ещё до
полного решения уравнений.

Если после безразмеризации возникает малый параметр
\[
\varepsilon A(u)+B(u)=0,
\qquad \varepsilon\ll1,
\]
можно попытаться построить предельную модель $B(u)=0$. Однако необходимо
проверять регулярность предельного перехода: в сингулярно возмущённых
задачах малый коэффициент может порождать пограничные слои и другие эффекты.

\subsection{Универсальные структуры}

Наиболее распространённая структура --- закон баланса:
\[
\boxed{
\frac{dQ}{dt}
=\text{приток}-\text{отток}+\text{источники}.
}
\]
В распределённой системе он принимает форму
\[
\frac{\partial u}{\partial t}
+\nabla\cdot\mathbf J=s.
\]
Из этой схемы получаются уравнения сохранения массы, энергии, заряда,
численности популяции и многие экономические балансовые модели.

Другие часто повторяющиеся структуры:
\begin{itemize}
    \item задачи минимизации энергии и вариационные постановки;
    \item линейные системы и собственные значения;
    \item диффузионные и волновые уравнения;
    \item графы и сетевые потоки;
    \item марковские и другие вероятностные модели.
\end{itemize}
Благодаря этому один численный метод, например метод конечных элементов,
может применяться к упругости, теплопроводности, электромагнитным и другим
задачам после изменения физической интерпретации коэффициентов.

\subsection{Перенос модели и границы универсальности}

Перенос модели в новую предметную область требует
\begin{enumerate}
    \item установить соответствие переменных и параметров;
    \item проверить размерности и знаки коэффициентов;
    \item заново сформулировать начальные и граничные условия;
    \item определить область применимости допущений;
    \item провести независимую валидацию.
\end{enumerate}

Например, логистическое уравнение
\[
\dot N=rN\left(1-\frac NK\right)
\]
может описывать рост популяции или другой процесс с насыщением, но значения
$r$ и $K$ приобретают совершенно иной предметный смысл. Совпадение формулы
не освобождает от проверки адекватности.


\subsection{Построение безразмерных комплексов через матрицу размерностей}

Пусть величины $q_1,\ldots,q_n$ выражаются через $r$ независимых основных
размерностей. Запишем их размерности в виде
\[
    [q_j]=D_1^{a_{1j}}\cdots D_r^{a_{rj}}
\]
и составим матрицу показателей
\[
    A=(a_{ij})\in\mathbb R^{r\times n}.
\]
Произведение
\[
    \Pi=q_1^{\alpha_1}\cdots q_n^{\alpha_n}
\]
безразмерно тогда и только тогда, когда
\[
    \boxed{A\alpha=0.}
\]
Следовательно, число независимых безразмерных произведений равно
\[
    \dim\ker A=n-\operatorname{rank}A=n-r_A.
\]
Это линейно-алгебраическая основа $\Pi$-теоремы Бэкингема и практический способ
строить критерии подобия.

\subsection{Геометрическое, кинематическое и динамическое подобие}

Для переноса результата с модели на натурный объект обычно требуется несколько
уровней подобия:
\begin{itemize}
    \item \textbf{геометрическое} --- одинаковая безразмерная форма и отношения
    характерных размеров;
    \item \textbf{кинематическое} --- подобие полей скоростей и временных
    масштабов;
    \item \textbf{динамическое} --- одинаковые отношения существенных сил,
    то есть совпадение соответствующих безразмерных критериев.
\end{itemize}
Например, для вязкого течения совпадение $Re$ необходимо для динамического
подобия по отношению инерционных и вязких сил; при существенной свободной
поверхности дополнительно возникает число Фруда. Нельзя автоматически сохранить
все критерии при изменении масштаба, поэтому физическое моделирование часто
требует выбора доминирующих механизмов.

\subsection{Безразмеризация как способ построения предельных моделей}

После безразмеризации уравнения часто принимают вид
\[
    A(u)+\varepsilon B(u)=0.
\]
Если $\varepsilon\ll1$ и предельный переход регулярный, первым приближением
служит более простая модель
\[
    A(u)=0.
\]
Так возникают иерархии моделей: несжимаемый предел при малом числе Маха,
невязкий предел при большом числе Рейнольдса вдали от стенок, квазистатические
приближения при малой роли инерции и т.п. Однако предельный переход может быть
сингулярным: например, малая вязкость порождает пограничный слой. Поэтому
универсальная предельная модель нуждается в проверке равномерности приближения.

\subsection{Универсальные канонические структуры}

Повторяемость математических моделей проявляется не только в совпадении одного
уравнения. Наиболее типичны несколько классов:
\begin{itemize}
    \item \textbf{законы баланса}
    $\partial_t u+\nabla\cdot J=s$;
    \item \textbf{градиентные и вариационные системы}, где движение или
    равновесие определяется энергией;
    \item \textbf{линейные спектральные задачи}
    $Au=\lambda Bu$;
    \item \textbf{диффузионные уравнения}, сглаживающие пространственные
    неоднородности;
    \item \textbf{волновые и гиперболические системы} с конечной скоростью
    распространения;
    \item \textbf{сетевые модели} и законы сохранения потоков на графах;
    \item \textbf{марковские модели}, задающие эволюцию распределений
    вероятностей.
\end{itemize}
Поэтому методы анализа и численные алгоритмы часто переносятся между физикой,
биологией, экономикой и инженерией при новой интерпретации переменных.

\subsection{Универсальность и границы переноса параметров}

Даже если форма уравнений совпадает, численные значения коэффициентов и
структура допустимых режимов могут быть совершенно различны. Универсальность
математической формы не означает универсальности параметров. При переносе
необходимо заново проверять:
\begin{itemize}
    \item размерности и безразмерные критерии;
    \item знак и допустимый диапазон коэффициентов;
    \item устойчивость и характерные времена;
    \item начальные и граничные условия;
    \item идентифицируемость параметров по данным новой предметной области.
\end{itemize}
Именно поэтому теория подобия облегчает перенос результатов, но не отменяет
валидацию новой интерпретации модели.

\subsection[Короткая схема доказательной идеи Пи-теоремы]{Короткая схема доказательной идеи $\Pi$-теоремы}

Размерностно однородная физическая зависимость не должна зависеть от выбора
единиц измерения. Изменение основных единиц действует на размерные величины как
$r$-параметрическая группа масштабных преобразований. Инвариантами этого действия
являются как раз мономы $q^\alpha$ с $A\alpha=0$. Пространство независимых
инвариантов имеет размерность $n-r_A$, поэтому исходную зависимость можно
выразить через $n-r_A$ независимых безразмерных комбинаций. Это объясняет, почему
теорема уменьшает число существенных аргументов без решения самих уравнений
процесса.

\subsection{Что важно сказать на экзамене}

Нужно объяснить, что универсальность возникает из абстрагирования и
повторяемости математических структур. Центральный строгий результат ---
$\Pi$-теорема Бэкингема. Полезно привести один пример одинакового уравнения
в разных областях, один пример безразмерного критерия и обязательно
подчеркнуть, что перенос модели требует повторной проверки области
применимости и валидации.

\newpage
